%% kapa: exam=1 round=24 year=2013 period=2 form=A subject=accounting subject_ko=회계학 seq=40 number=40
%% source: Q-Net 감정평가사 기출문제 (cst00309) artlSeq=5073802 file=extracted/1008/5073802/2013년 제24회 감정평가사 1차시험 기출문제/2교시(A형).hwp
%% conversion: kordoc 4.24.51 -> tools/gen1.py (auto); review=proofread-sonnet; flags=none
\begin{question}{40}
(주)감평은 품질관련 활동원가를 예방원가, 평가원가, 내부실패원가 및 외부실패원가로 구분하고 있다. 다음에 제시한 자료 중 외부실패원가로 집계할 금액은?
\begin{framedtext}
\begin{tabular}{@{}l@{\hspace{4em}}r@{}}
제품보증수리활동 & ￦21,000 \\
원재료 검사활동 & 11,000 \\
직원 품질교육활동 & 50,000 \\
고객서비스센터활동 & 6,000 \\
불량품 재작업활동 & 8,000 \\
설비보수 및 유지활동 & 5,000 \\
판매기회 상실로 인한 기회비용 & 18,000 \\
공정검사활동 & 7,000 \\
설계개선활동 & 10,000
\end{tabular}
\end{framedtext}

\begin{choices}
  \item ￦35,000
  \item ￦43,000
  \item ￦45,000
  \item ￦53,000
  \item ￦57,000
\end{choices}
\end{question}
\answer{3}
